\documentclass[runningheads]{llncs}
\usepackage{graphicx}
\usepackage{hyperref}
\begin{document}
\title{A Clean Test Paper}
\author{Jane Doe}
\institute{Some University}
\maketitle
\begin{abstract}
This abstract describes scenario-based testing of automated vehicles in a few words.
\end{abstract}

\section{Introduction}\label{sec:intro}
Automated vehicles are complex systems that must be tested thoroughly before release.
Scenario-based testing (SBT) is a common approach for this purpose.
In this paper, we apply SBT to a new class of driving functions and report on three case studies.
Figure~\ref{fig:arch} gives an overview of the architecture that we use throughout the paper.
The evaluation in Sect.~\ref{sec:eval} shows that the approach scales to realistic traffic situations.
We use SBT with several tools (e.g.,~\cite{smith2020,doe2021}) and compare the results with
related approaches (cf.~\cite{lee2019}).
The tool is available online.\footnote{\url{https://zenodo.org/records/1234567}, last accessed 2026-01-15.}

The remainder of this paper is structured as follows: Sect.~\ref{sec:related} discusses related
work, Sect.~\ref{sec:method} presents our method, Sect.~\ref{sec:eval} evaluates it, and
Sect.~\ref{sec:conc} concludes the paper.

\section{Related Work}\label{sec:related}
Smith~\cite{smith2020} introduced a framework for testing automated vehicles that relies on
a catalog of traffic situations. Doe~\cite{doe2021} extended this framework with a formal
description language for scenarios. Both approaches focus on the functional behavior of the
vehicle and do not consider the interaction between several driving functions, which is the
gap that our work addresses in the following sections of this paper.

\section{Method}\label{sec:method}
Our method consists of three steps that are executed for every driving function under test.
First, we derive relevant scenarios from the requirements of the function. Second, we execute
these scenarios in a simulation environment and record the behavior of the vehicle. Third, we
evaluate the recorded behavior against the requirements and report every violation that occurs.
\begin{figure}[t]
\centering
\includegraphics[width=\textwidth]{arch}
\caption{Architecture of the test framework.}\label{fig:arch}
\end{figure}

\section{Evaluation}\label{sec:eval}
We evaluate the method on three driving functions and report the results in Table~\ref{tab:res}.
The method detected every violation that was known beforehand and found two new ones, which the
developers of the driving functions confirmed after a manual inspection of the recorded runs.
\begin{table}[t]
\centering
\caption{Results of the evaluation.}\label{tab:res}
\begin{tabular}{lr}
Function & Violations \\
A & 3 \\
\end{tabular}
\end{table}

\section{Conclusion}\label{sec:conc}
We presented a method for SBT of interacting driving functions and showed in
three case studies that it detects known and unknown violations. In future work, we plan to apply
the method to further driving functions and to integrate it into a continuous integration setup.

\bibliographystyle{splncs04}
\bibliography{refs}
\end{document}
